% ---------------------------------------------------------------------------
% The forward direction in three pages: grammars, the transpiler, coverage, and
% the emptying of the inexpressible column. Tables are the complete paper's.
% ---------------------------------------------------------------------------

\section{The Forward Direction, Briefly}
\label{s:sec:forward}

The forward direction is reported compactly here because its answer is negative in the
useful sense: it found no obstruction. This section says what was built, what it covers,
and why the shortfall is not evidence of inexpressibility --- which is what makes
\S\ref{s:sec:regular} the interesting half.

\paragraph{The grammars exist, and they are deterministic.}
Neither vendor publishes a grammar. We infer one for each language, and the Datadog
grammar is, as far as we have found, the first published for it --- which is a statement
about a search and is reported as one; the vendor documents the language only by example.
Both grammars admit deterministic one-token-lookahead parsers, at the sizes
Table~\ref{s:tab:grammars} reports. Datadog's is conflict-free. PromQL's carries two shift/reduce conflicts on a left parenthesis, both
resolved the way the reference server resolves them, and the complete paper documents the
resolution and shows it is the one the published implementation makes. That matters only
because it licenses what follows: a unique parse tree per query is what makes a
syntax-directed translation well defined at all.

\begin{table}[htbp]\centering\footnotesize
\begin{tabular}{@{}llrrrl@{}}
\toprule
Language & Grammar & Terminals & Rules & LALR states & Conflicts \\
\midrule
promql & v1 & 32 & 73 & 130 & 2$\times$ \texttt{Shift/Reduce on LPAR} \\
promql & v2 & 33 & 82 & 145 & 2$\times$ \texttt{Shift/Reduce on LPAR} \\
datadog & v1 & 19 & 35 & 66 & none \\
datadog & v2 & 22 & 60 & 111 & none \\
\bottomrule
\end{tabular}

\caption{The two inferred grammars, at the two revisions this work produced for each.
Conflicts are LALR(1) conflicts reported by the generator; Datadog's grammar has none at
either revision, and PromQL's two are the shift/reduce pair on a left parenthesis that the
reference server resolves the same way.}
\label{s:tab:grammars}
\end{table}

\paragraph{The translator.}
$T$ factors as $\pi_1\circ E\circ L$ through an intermediate representation, as
Definition~\ref{s:def:translator} sets out. The lowering reads a Datadog query into the
IR; the emission writes PromQL and returns, beside the expression, a set of
\emph{caveats} recording what it had to assume. Every refusal is a $\bot$ with a gap
code. The IR is the only object in this work whose denotation is written out, and the
exactness results of the complete paper are proved over $E(L(q))$ rather than over $T(q)$,
because they need the caveat set that $E$ returns alongside the expression.

\paragraph{One forward result the reverse direction needs.}
The scope syntax of Datadog admits a disjunction over a single tag key, written
$k \in (g_1,\dots,g_B)$ or $k{:}g_1 \vee \cdots \vee k{:}g_B$, and this is the one
forward-direction result \S\ref{s:sec:regular} uses, so it is stated here rather than
cited.

\begin{proposition}[Scope disjunction is one matcher]\label{prop:disjunction}
A Datadog scope disjoining globs $g_1,\dots,g_B$ over one tag key $k$ denotes exactly the
set denoted by the single PromQL matcher
$k\,\texttt{=\textasciitilde{}}\,\mathrm{glob}(g_1)\texttt{|}\cdots\texttt{|}\mathrm{glob}(g_B)$,
under PromQL's whole-value anchoring.
\end{proposition}

\noindent
The argument is one line in each direction: a value matches the scope exactly when it
matches some $g_i$, and matches the alternation exactly when it matches some
$\mathrm{glob}(g_i)$, which Proposition~\ref{prop:glob-regex} identifies with $L(g_i)$.
The anchoring is read off a live server rather than assumed, and the complete paper
reports that measurement. Read in the other direction the same proposition says that the
reverse translation of a regular expression should be sought as a \emph{finite union} of
globs --- the two directions are not two results but one fact used twice.

\paragraph{What it covers.}
The corpus is $22{,}894$ queries extracted from published dashboard and alerting
documents, of which $8{,}770$ are Datadog metric queries. Table~\ref{s:tab:coverage}
gives the outcome for every occurrence.

\begin{table}[htbp]\centering\footnotesize
\begin{tabular}{@{}lrrrr@{}}
\toprule
Outcome & Occurrences & Share & Distinct & Share \\
\midrule
translated ($\kappa$) & 8,376 & 95.6\% & 7,345 & 95.5\% \\
$\bot$ (gap catalogue) & 382 & 4.4\% & 337 & 4.4\% \\
unparsed by the §4 grammar & 8 & 0.1\% & 6 & 0.1\% \\
\bottomrule
\end{tabular}

\caption{Outcomes over the $8{,}770$ Datadog metric query occurrences. Shares are of
occurrences and of distinct query strings respectively; a query used on forty dashboards
counts forty times in the first and once in the second.}
\label{s:tab:coverage}
\end{table}

So $\kappa=96.6\%$. The grammar and the implementation order were both fitted to this
corpus, which makes that a resubstitution estimate and not a prediction. Frozen and run
on a disjoint repository of dashboards by other authors it returns $96.7\%$, refusing
only in ways the gap catalogue already names --- a tenth of a point from the figure here,
against $1.1$ points of dashboard-clustered standard error resting on an effective two
dashboards. Nothing large is hiding on that side, which is a weaker claim than the two
rates agreeing, and it is the one the data supports.

\paragraph{The whole shortfall is named, and none of it is inexpressibility.}
The $3.3\%$ that refuses is not a residue. It is five gap codes, listed in
Table~\ref{s:tab:gaps}, and each names a construct rather than a difficulty.

\begin{table}[htbp]\centering\scriptsize
\begin{tabular}{@{}lrrll@{}}
\toprule
Gap & Occurrences & Distinct & Share & Shortest example \\
\midrule
\texttt{UNSUPPORTED\_FUNCTION} & 312 & 270 & 3.6\% & \texttt{sum:iis.users.anon\{*,!*\}} \\
\texttt{FILL\_INTERPOLATION} & 36 & 36 & 0.4\% & \texttt{avg:mongodb.mem.resident\{*,*\}.fill(zero)} \\
\texttt{LABEL\_NAME\_GLOB} & 31 & 28 & 0.4\% & \texttt{sum:velero.backup.items\{*-location,*-name,*-u\dots{}} \\
\texttt{SCOPE\_DISJUNCTION} & 4 & 4 & 0.0\% & \texttt{sum:cert\_manager.http\_acme\_client.request.cou\dots{}} \\
\bottomrule
\end{tabular}

\caption{The complete gap catalogue for the forward direction. Five codes account for
every refusal. Shares are of the $8{,}770$ occurrences.}
\label{s:tab:gaps}
\end{table}

\noindent
This paper reserved a sixth column for constructs beyond PromQL's reach, and over the
course of the work that column emptied of everything we can demonstrate, bar $14$
occurrences left undecided. Each candidate for it was removed in one of two
ways: an implementation that produced an exact PromQL expression for the construct, which
shows the refusal was never a gap; or an argument that the construct is expressible,
which shows the same thing without the code being written. What survives is $14$
occurrences left undecided --- a quantile construct whose Datadog semantics we could not
pin down well enough to say either way --- and they are reported as undecided rather than
assigned to either side. Remark~\ref{s:rem:refusal} is the reason this distinction is
drawn so insistently: before those removals, each of those constructs had a gap code, and
a gap code reads like a fact about the languages.

It is worth being exact about what that does and does not establish. It does not say
PromQL can express every Datadog metric query; it says that among the constructs this
corpus contains, we can demonstrate no case where it cannot, and that the demonstrations
run in the direction of expressiveness rather than away from it. Most of the corrections
that emptied the column were to our own earlier claims rather than to the vendor's
behaviour, which is the kind of finding that is easy to leave unsaid.

A translator's willingness is also not the same as its correctness, and the complete
paper separates the two: it runs every distinct emitted query shape against a live
Datadog account and a local Prometheus and compares results under
Definition~\ref{s:def:d}. That run found four translator defects, and once they were
fixed the willingness measure and the agreement measure coincided on every shape it could
compare. The forward number above is therefore a coverage figure with a correctness check
behind it, which is the strongest form in which this work can offer it.
